\documentclass[11pt]{article}
\usepackage{preamble}
\renewcommand{\answer}[1]{}

\begin{document}

\ladderheading{AP Statistics \textperiodcentered\ Unit 4 \textperiodcentered\ Topic 4.2 \textperiodcentered\ B: 2027-01-12 \textperiodcentered\ E: 2027-02-26}{Can Chance Explain the Mean Difference?}

\focusbanner{TODAY'S PROBLEM}

Suppose a researcher recruits 50 student volunteers to test whether a planning prompt lowers mean logic-task time, then randomly assigns 25 to the prompt and 25 to standard directions. Their means are 14.2 and 17.8 minutes, so prompt minus standard is $-3.6$. Under no effect, 18 of 2{,}000 rerandomized differences are $-3.6$ or lower.\par\smallskip \textbf{Nia:} ``A result this favorable is rare under no effect, supporting a causal decrease for students like these volunteers.''\par \textbf{Owen:} ``The $0.009$ is the probability of no effect, so there is a $99.1\%$ chance the prompt works for every district student.''\par
\begin{center}
\textbf{Experiment results}\par\smallskip
\begin{tabular}{|l|c|c|}
\hline
\textbf{Group} & \textbf{n} & \textbf{Mean (min)} \\ \hline
\textbf{Prompt} & 25 & 14.2 \\ \hline
\textbf{Standard} & 25 & 17.8 \\ \hline
\end{tabular}
\end{center}
\begin{firsttakebox}{FIRST TAKE --- before the video (5 min)}
Whose use of the rerandomization result is better supported, Nia's or Owen's? Cite $-3.6$ or $18/2{,}000$.
\workspace{0.9in}
\end{firsttakebox}

\begin{callout}{\IconBook\ ASK WHAT CHANCE WOULD DO UNDER NO EFFECT (keep on the page)}{calloutgreen}
In a randomized experiment, the observed difference in sample means is evidence for a directional claim.
Under no treatment effect, rerandomize the recorded responses into groups of the original sizes, compute the
difference in means in the stated order, and repeat. The proportion at least as extreme in the claim's direction
estimates $P(\text{result this extreme}\mid\text{no effect})$; it is not the probability that no effect is true.
A small tail proportion weakens the chance-only explanation but cannot make error impossible. Random assignment
supports causation for the experimental units; random sampling is needed to generalize to a target population.
\end{callout}

\newpage
\textbf{Finish after the video:}\par\smallskip

\dokbadge{1}~\textbf{(a)} State the directional claim, the observed evidence statistic, and its expected value under no effect.\par
\workspace{0.7in}

\dokbadge{2}~\textbf{(b)} Describe one rerandomization trial, compute the tail proportion, and interpret it under no effect.\par
\workspace{1.4in}

\dokbadge{3}~\textbf{(c)} Decide whose conclusion is warranted using 0, $-3.6$, and 18 of 2{,}000. Distinguish the two directions of conditional probability, bound causation and generalization from the design, identify the remaining inference error, and explain the rejected statement's consequence.\par
\workspace{2.0in}

\begin{sentenceframebox}\raggedright\textbf{Frame:} The $0.009$ estimates the chance of \blank[2.2in] assuming \blank[1.8in]; it does not estimate \blank[2.0in].\end{sentenceframebox}

\begin{sentenceframebox}\raggedright\textbf{Frame:} Random assignment supports \blank[1.7in], but volunteer recruitment prevents \blank[2.2in].\end{sentenceframebox}

\begin{summaryexitbox}{EXIT --- turn this in}
One thing the video changed about my first take: \hrulefill\par
\workspace{0.4in}
\end{summaryexitbox}

\end{document}
